% GERADO por tools/generation/gerar_secoes_latex.py a partir de
% artifacts/paper/secoes/07-avaliacao.md. Nao editar a mao: a proxima geracao
% desfaz. Corrija o Markdown e regere.
\subsection{As questões de competência sobre o cenário instanciado}

A instanciação deriva cada indivíduo de um trecho identificado da publicação do cenário \cite{sbsi_estendido}: o observatório e seu gerenciamento de conteúdo, as oito funcionalidades, os três perfis de usuário levantados com 25 participantes e o sistema institucional que não se integra a ele. O que foi acrescentado está declarado item a item --- uma segunda fonte de dados, a cadeia de ETL com carimbos de tempo, um segundo componente de gerenciamento, dois projetos-exemplo e duas observações datadas ---, e o mapa entre trecho e indivíduo acompanha o depósito, junto do observatório municipal de obras públicas que a comparação da QC7 exigiu, sintético e declarado como limitação (\S{}8). As sete consultas retornaram resultado não-vazio e semanticamente correto (Tabela 3), e nenhuma interroga a arquitetura do artefato.

\begin{table}[H]
\centering
\caption{As sete questões de competência sobre o cenário instanciado.}
\label{tab:3}
\scriptsize
\begin{tabular}{>{\centering\arraybackslash}p{\dimexpr0.0706\linewidth-2\tabcolsep\relax}p{\dimexpr0.4588\linewidth-2\tabcolsep\relax}>{\centering\arraybackslash}p{\dimexpr0.0706\linewidth-2\tabcolsep\relax}p{\dimexpr0.4000\linewidth-2\tabcolsep\relax}}
\hline
 & \textbf{Pergunta} & \textbf{Linhas} & \textbf{Construto de que depende} \\
\hline
QC1 & Que interfaces um gerenciamento provê? & 8 & \texttt{ViewProvision}, de A9 \\
QC2 & Quem acessa que conteúdo de um projeto? & 4 & \texttt{StakeHolder} «roleMixin», de A6 \\
QC3 & Quem carregou uma fonte, e quando? & 4 & o ETL como evento, de A5 \\
QC4 & Que projetos ficaram sem observação? & 1 & \texttt{Observation} como evento, de B8 \\
QC5 & Que motivações movem cada tipo de ator? & 7 & a taxonomia de \texttt{Agent}, de A6 e A7 \\
QC6 & Qual a proveniência de um conteúdo? & 4 & a parthood de \texttt{EtlProcess}, de A5 \\
QC7 & Dois observatórios cobrem o mesmo? & 31 & as classes revisadas, com a C2 \\
\hline
\end{tabular}
\end{table}

As consultas e seus resultados completos estão no Apêndice A e, na íntegra, no depósito aberto, com um script que as reexecuta e confere cada uma contra o arquivo gravado.

\textbf{As divergências foram registradas, inclusive as que desfavorecem o artefato.} Cada questão correu contra expectativa escrita antes, e quatro apontam limite do modelo revisado: a QC2 responde sobre o recorte do relator que medeia agente e interface de relacionamento, não sobre acesso amplo; na QC3, quem participa da carga são papéis de um componente de software, não pessoa; a QC4 só distingue um projeto do outro porque o cenário dá a cada um um gerenciamento diferente; e a cadeia da QC6 abre em leque, entregando todos os processos de ETL do gerenciamento responsável, não o que produziu o conteúdo.

A QC7 sustenta a tese: dos 31 conceitos que ao menos uma das iniciativas instancia, \textbf{23 são comuns e oito divergem} --- conhecimento, grupo do observatório, interação social, observação, parte interessada e relacionamento aparecem só no primeiro; agente social e organização, só no segundo. Duas iniciativas aderentes ao mesmo modelo cobrem recortes conceituais distintos, e a divergência deixou de ser argumento para virar medida.

\subsection{Conformidade à UFO, antes e depois}

O verificador de conformidade sintática e semântica devolve \textbf{zero problema sobre os três modelos}, inclusive sobre a linha de base defeituosa, e apresentar esse zero como efeito da revisão seria enganoso: nenhuma das suas 24 regras alcança restrição meronímica, aridade de relator ou a distinção entre endurante e perdurante. O que ele atesta é que a revisão trocou dois estereótipos por camada inteira sem quebrar nenhuma das regras que ele tem; quatro mutações conhecidas mostram que ele acusaria se elas quebrassem.

A distância entre os modelos é medida por nove regras estruturais escritas para esta análise, que operacionalizam restrições fora daquele alcance e correm sobre os três modelos pelo mesmo caminho de código, e cai de \textbf{12 na linha de base para 3 depois da primeira rodada e 0 depois da segunda}: somem as relações de membro com o coletivo na ponta errada, a meronímica sem extremidade declarada, as mediações com mínimo zero, a colisão de nome com o metaconceito, os relatores acima de binário e as duas regras da ausência da microteoria de eventos. Sete delas medem, e têm por controle positivo a própria linha de base --- a execução reprova se alguma deixar de disparar ali; as outras duas vigiam um evento sem participante e uma participação invertida, que só a revisão produz, e seu controle é uma mutação do revisado, acusada nas duas.

\subsection{\textit{Pitfalls} na OWL, antes e depois}

As duas execuções do detector de \textit{pitfalls} \cite{poveda2014oops} submetem a mesma coisa pelo mesmo caminho de código, e é essa igualdade que os torna comparáveis. O total de elementos apontados cai de \textbf{105 para 56}, e a leitura do artefato de domínio é mais nítida: entre os elementos que o serviço lista, os que pertencem à OntoMPO caem de \textbf{60 para zero} --- eram todos ausências de anotação, fechadas pela C2. Dois \textit{pitfalls} ficam fora dessa contagem porque o serviço não lista os elementos afetados: a convenção de nome, deliberadamente não fechada, e a ausência de inversa declarada, que cai de 21 para 7 --- o número de propriedades emprestadas da gUFO que o arquivo declara.

A diferença entre as leituras não é ruído, e o custo precisa estar escrito: ao declarar localmente os termos da ontologia de fundamentação, a C1 os transforma em elementos que o serviço passa a avaliar --- sem anotação própria e, num caso, sem domínio nomeado. Fechar isso exigiria embutir a gUFO no arquivo, que é o que a importação existe para evitar. Pela mesma razão, o \textit{pitfall} de classe não tipada sobe de 13 para 16: some a metade nomeada e sobra a anônima. Por fim, os dois artefatos OWL são \textbf{consistentes}: um raciocinador do perfil OWL 2 RL não encontra contradição em nenhum deles, e oito mutações confirmam que ele acusaria --- quatro nos dois e quatro só no customizado, medida do conteúdo lógico que a customização acrescentou (\S{}8).
